\begin{description}
\item[(OP01.1)]
The orbits of the 4 TNOs are elliptical, which can be noticed from careful observations of
their speed. Hence, to estimate lengths of their semi-major axes it is essential to measure
the time period of one complete revolution of each TNO. The time periods of the planets are
\[
\mathrm{O1{:}\ 40\ sec},\quad \mathrm{O2{:}\ 30\ sec},\quad
\mathrm{O3{:}\ 48\ sec},\quad \mathrm{O4{:}\ 34\ sec}
\]
Hence, the correct ascending order is: $\mathrm{O2} < \mathrm{O4} < \mathrm{O1} < \mathrm{O3}$.
% The 24-row credit table for all possible orderings (marking scheme) is omitted.

\item[(OP01.2)]
The orbits of the TNOs can be drawn by marking a few important constellations and noticing
the fact that some prominent or important stars lie along the TNOs' paths. Also the
intersection points of each TNO trajectory with respect to the Earth's orbit (i.e.\ the
Ecliptic) lie near prominent stars.
% FIGURE NEEDED: traced TNO orbits on the sky map with tolerance bands
% (2025_Map-OP01-Solutions.pdf, overlay of Map-OP01).

Since this map has an angle preserving (Stereoscopic) projection, measuring the angles
directly using a protractor at the intersection points gives the inclination of the orbits
with respect to the Ecliptic.

The measured values of inclination are:
\begin{itemize}
\item $i_1 = 18^\circ$
\item $i_2 = 14^\circ$
\item $i_3 = 20^\circ$
\item $i_4 = 12^\circ$
\end{itemize}
The average inclination $i_{hp} = 16^\circ$.
% The per-TNO trace/angle marking scheme lines are omitted.
\end{description}
